\begin{problem}{G}{Gambler's Trick}{1 second}

\noindent
Kaito and Aoko are playing a board game to decide who will pay for their meal at the cafe.

\noindent
The rules are simple. There are $M$ squares numbered from $1$ to $M$, and a $D$-sided die whose faces are numbered from $1$ to $D$. Each player controls a character, and they play independently.

\noindent
Initially, both characters are on square $A$. On each turn, a player rolls the die. If the result is $d$ and their character is currently on square $x$, the character moves to square $x+d$. If $x+d \ge M$, the player \textbf{immediately loses} and their game ends. Otherwise, the player \textbf{must} continue playing until they have completed $N$ turns.

\noindent
Of course, Kaito has a trick up his sleeve. He can secretly control the die and make it show any value from $1$ to $D$.

\noindent
Let $P_i$ be Kaito's position after the $i$-th turn. Kaito wants to show off by moving as far ahead as possible as early as possible. Therefore, among all ways for him to survive all $N$ turns, he wants the sequence $P_1,P_2,\ldots,P_N$ to be \textbf{lexicographically largest}.

\noindent
For two different arrays $X$ and $Y$ of the same length, $X$ is \textbf{lexicographically larger} than $Y$ if, at the first position where they differ, $X$ has the larger value.

\noindent
Find Kaito's desired sequence, or print $-1$ if it is impossible for him to survive all $N$ turns.

\InputFile

\noindent
The first line contains an integer $T$, the number of test cases.

\noindent
Each test case consists of a single line containing four integers $N$, $D$, $A$, and $M$ ($1 \le N \le 2 \cdot 10^5,\ 4 \le D \le 10^9,\ 1 \le A < M \le 10^9$) --- the number of turns, the number of faces on the die, the initial square of the characters, and the losing square, respectively.

\noindent
It is guaranteed that the sum of $N$ over all test cases does not exceed $2 \cdot 10^5$.

\OutputFile

\noindent
For each test case, print one line.

\noindent
If it is impossible for Kaito to survive all $N$ turns, print $-1$.

\noindent
Otherwise, print $N$ integers $P_1, P_2, \ldots, P_N$ --- Kaito's lexicographically largest possible sequence of positions.

\Examples

\begin{example}
\exmp{1}{5
3 4 1 14
3 4 5 14
3 4 10 14
3 4 11 14
3 4 2 16
}{5 9 13
9 12 13
11 12 13
-1
6 10 14
}
\end{example}

\pagebreak

\end{problem}
